\documentclass[11pt]{article}
\usepackage[margin=1in]{geometry}
\usepackage{amsmath,amssymb,amsthm}
\usepackage[hidelinks]{hyperref}
\newtheorem{theorem}{Theorem}
\newtheorem{corollary}[theorem]{Corollary}
\newcommand{\bin}[2]{\binom{#1}{#2}}
\title{Proof of Mathar's conjecture on OEIS A027182:\\
a closed form and the six-term recurrence}
\author{Haobo Ma\thanks{ChronoAI Pte Ltd, Singapore.}\and
Wenlin Zhang\thanks{National University of Singapore; ChronoAI Pte Ltd, Singapore.}}
\date{1 November 2026}
\hypersetup{pdftitle={Proof of Mathar's conjecture on OEIS A027182: a closed form and the six-term recurrence},
pdfauthor={Haobo Ma; Wenlin Zhang},pdfsubject={CC BY 4.0}}
\begin{document}
\maketitle
\begin{abstract}
We prove the recurrence for A027182 stated as a conjecture by R.\ J.\ Mathar
(OEIS A027182, FORMULA section, 6 May 2026) for all $n\geq6$, by evaluating
the middle coefficient of the square of each A027170 row polynomial.
Row symmetry reduces the convolution to five binomial sums, and a shifted
second-order recurrence gives the six-term recurrence. The same
central-binomial identity proves Mathar's conjectured recurrences for
A027171, A027174 and A027175; a parity decomposition also proves his
seven-term recurrence for the half-row sums A027179.
\end{abstract}

Clark Kimberling's sequence A027182 is defined from the row polynomials
of A027170. Mathar conjectured a polynomial-coefficient recurrence for
it on 6 May 2026~\cite{A027182}. We prove that recurrence by first
evaluating the convolution explicitly. General algorithms for deriving
recurrences of binomial sums are developed by Bostan, Lairez and
Salvy~\cite[Introduction and \S4]{BLS}; here the row symmetry permits
a direct evaluation by the binomial theorem.
The related closed forms of Christian Krause in A027171, A027174 and
A027175~\cite{A027171,A027174,A027175} reduce three further conjectures of
Mathar, also dated 6 May 2026, to the same central-binomial identity.
The half-row conjecture in A027179~\cite{A027179} follows after separating
even and odd rows. We give those deductions below, crediting the recorded
closed forms separately from the recurrence proofs.

For integers $n\geq0$ and $0\leq k\leq n$, follow A027170~\cite{A027170}
and set
\begin{align*}
t(n,k)&=\bin{n+2}{k+1}+\bin{n+1}{k}+\bin{n+1}{k+1}+\bin nk,\\
T(n,k)&=t(n,k)-t(0,0)+1=t(n,k)-4,\\
g_n(x)&=\sum_{k=0}^n T(n,k)x^k,\qquad a_n=[x^n]g_n(x)^2.
\end{align*}
Here $[x^n]$ denotes coefficient extraction. The original A027182
indexing starts at $n=1$; the same definition gives $a_0=1$.

\begin{theorem}\label{thm:main}
For every integer $n\geq0$,
\begin{equation}\label{eq:closed}
a_n=4\bin{2n+4}{n+2}+4\bin{2n+2}{n+1}+\bin{2n}{n}
       -72\cdot2^n+16n+40.
\end{equation}
Consequently, for every integer $n\geq6$,
\begin{equation}\label{eq:rec}
\begin{aligned}
0={}&2(n+2)a_n-(17n+23)a_{n-1}+(47n+36)a_{n-2}\\
 &-2(22n+17)a_{n-3}+4(20-n)a_{n-4}\\
 &+8(2n-9)a_{n-5}+24(12n-47).
\end{aligned}
\end{equation}
\end{theorem}

\begin{proof}
Put $A_k=\bin{n+2}{k+1}$ and $B_k=\bin nk$. Pascal's identity gives
$T(n,k)=2A_k+B_k-4$. Since $A_{n-k}=A_k$ and $B_{n-k}=B_k$,
the row is symmetric. Hence
\[
a_n=\sum_{k=0}^n T(n,k)T(n,n-k)=\sum_{k=0}^n(2A_k+B_k-4)^2.
\]
The required sums, all with $0\leq k\leq n$, are
\begin{align*}
\sum A_k^2&=\bin{2n+4}{n+2}-2,&
\sum A_kB_k&=\bin{2n+2}{n+1},&
\sum B_k^2&=\bin{2n}{n},\\
\sum A_k&=2^{n+2}-2,&\sum B_k&=2^n.
\end{align*}
For completeness, coefficient extraction from
$(1+x)^r(1+x)^s=(1+x)^{r+s}$ gives
$\sum_j\bin rj\bin s{q-j}=\bin{r+s}q$.
Taking $(r,s,q)=(n+2,n+2,n+2)$ gives the first squared sum after
removing its two endpoint terms. Taking $(r,s,q)=(n+2,n,n+1)$
and using $\bin nk=\bin n{n-k}$ gives the mixed sum; its omitted
$j=0$ term is zero. Taking $(r,s,q)=(n,n,n)$ gives the remaining
squared sum. The two ordinary sums follow by putting $x=1$ in the
binomial theorem, with two endpoint terms removed for $A_k$.
Expanding the square now gives
\[
a_n=4\sum A_k^2+4\sum A_kB_k+\sum B_k^2
       -16\sum A_k-8\sum B_k+16(n+1),
\]
which is exactly~\eqref{eq:closed}.

We next derive~\eqref{eq:rec}, keeping the variable coefficients inside
each shift. Write
\[
c_m=\bin{2m}m,\qquad h_m=4c_{m+2}+4c_{m+1}+c_m,
\qquad p_m=16m+40.
\]
For any sequence $f$, define
\[
(Bf)_m=2(m+4)f_{m+2}-(7m+27)f_{m+1}-(4m+2)f_m.
\]
The factorial definition of $c_m$ implies
$R_m:=(m+1)c_{m+1}-2(2m+1)c_m=0$.
Expanding coefficients of $c_m,\ldots,c_{m+4}$ yields
\[
(Bh)_m=8R_{m+3}+12R_{m+2}+6R_{m+1}+R_m=0.
\]
Let $(Ef)_m=f_{m+1}$ and $Q(E)=E^3-5E^2+8E-4=(E-1)(E-2)^2$.
Define $L=Q(E)B$, meaning that $E$ shifts the entire sequence $Bf$.
Expansion gives
\begin{equation}\label{eq:L}
\begin{aligned}
(Lf)_m={}&2(m+7)f_{m+5}-(17m+108)f_{m+4}\\
 &+(47m+271)f_{m+3}-(44m+254)f_{m+2}\\
 &+(60-4m)f_{m+1}+(16m+8)f_m.
\end{aligned}
\end{equation}
In particular $Lh=0$. Direct substitution gives
\[
(B(2^{\bullet}))_m=(-10m-24)2^m,\qquad
(Bp)_m=-144m^2-744m-1016.
\]
The factor $(E-2)^2$ kills every sequence $(\alpha m+\beta)2^m$.
For a quadratic $u_m=\alpha m^2+\beta m+\gamma$, expansion gives
$(Q(E)u)_m=2\alpha m+\beta-3\alpha$.
Consequently,
\[
L(2^{\bullet})=0,\qquad (Lp)_m=-288m-312=-24(12m+13).
\]
Equation~\eqref{eq:closed} says $a_m=h_m-72\cdot2^m+p_m$.
Linearity therefore gives $(La)_m+24(12m+13)=0$ for every $m\geq0$.
Putting $m=n-5$ in~\eqref{eq:L} proves~\eqref{eq:rec} for $n\geq6$.
\end{proof}

\paragraph{Initial values and source convention.}
The definition gives
\[
(a_0,a_1,a_2,a_3,a_4)=(1,18,150,820,3726).
\]
The example in A027182 prints a final value of $120$ for
$5^2+10^2+5^2$; the sum is $150$, agreeing with the entry's term list
and with~\eqref{eq:closed}. With the extension $a_0=1$, the same
recurrence also holds at $n=5$.

\section*{Four companion conjectures}

The following corollaries prove the statements labeled ``Conjecture'' in
the FORMULA sections of A027171, A027174, A027175 and A027179, each
attributed there to R.\ J.\ Mathar on 6 May 2026.

\begin{corollary}\label{cor:middle}
Let $u_n=T(2n,n)$ for $n\geq0$ (A027171), and
$v_n=T(2n-1,n-1)$ for $n\geq1$ (A027175). Then
\begin{equation}\label{eq:middle}
2(n+1)f_n-(7n+5)f_{n-1}+2(3-2n)f_{n-2}+12(1-3n)=0
\end{equation}
holds for $f=u$, $n\geq2$, and for $f=v$, $n\geq3$.
\end{corollary}

\begin{proof}
Krause's formulas~\cite{A027171,A027175}, dated 10 and 11 May 2026,
respectively, follow immediately from the definition of $T$ and give
\[
u_n=d_n-4,\qquad v_n=\tfrac12d_n-4,\qquad
d_n=2c_{n+1}+c_n=\frac{9n+5}{n+1}c_n.
\]
For the second formula, use
$\bin{2n+1}{n}=(2n+1)c_n/(n+1)$ and
$\bin{2n-1}{n-1}=c_n/2$.
In terms of the already established $R_m=0$, expansion gives
\[
2(n+1)d_n-(7n+5)d_{n-1}+2(3-2n)d_{n-2}
=4R_n+4R_{n-1}+R_{n-2}=0.
\]
The sum of the three coefficients on the left is $3-9n$.
Subtracting $4$ therefore contributes $36n-12$, which is canceled by
$12(1-3n)$. The smallest indices in the statement keep all original
sequence indices in their domains.
\end{proof}

\begin{corollary}\label{cor:near}
For $z_n=T(2n,n-2)$, $n\geq2$ (A027174), and every $n\geq4$,
\begin{equation}\label{eq:near}
\begin{aligned}
0={}&-2(n+3)(139n-175)z_n
 +(1279n^2+1334n-1089)z_{n-1}\\
 &-2(167n+121)(2n-3)z_{n-2}+12(111n^2+456n+229).
\end{aligned}
\end{equation}
\end{corollary}

\begin{proof}
Krause's formula of 11 May 2026~\cite{A027174}, also immediate from $T$,
can be written
\[
z_n=2\bin{2n+2}{n-1}+\bin{2n}{n-2}-4=r(n)c_n-4,
\qquad r(n)=\frac{n(9n^2+14n+1)}{(n+1)(n+2)(n+3)}.
\]
Indeed, the ratios of the binomial coefficients themselves to $c_n$ are
$2n(2n+1)/((n+2)(n+3))$ and $n(n-1)/((n+1)(n+2))$;
the first has multiplier $2$ in $z_n$.
Denote the three coefficients of $z_n,z_{n-1},z_{n-2}$ in
\eqref{eq:near} by $q_0,q_1,q_2$. Since
\[
\frac{c_{n-1}}{c_n}=\frac{n}{2(2n-1)},\qquad
\frac{c_{n-2}}{c_n}=\frac{n(n-1)}{4(2n-1)(2n-3)},
\]
the homogeneous part for $r(n)c_n$ vanishes by the polynomial identity
\begin{equation}\label{eq:nearcert}
\begin{aligned}
0={}&-4n(2n-1)(139n-175)(9n^2+14n+1)\\
 &+(n-1)(9n^2-4n-4)(1279n^2+1334n-1089)\\
 &-(n-2)(n+2)(167n+121)(9n^2-22n+9).
\end{aligned}
\end{equation}
More precisely, its left side is
$2(n+1)(n+2)(2n-1)\sum_{j=0}^2q_j(n)r(n-j)c_{n-j}/c_n$;
expanding the displayed products gives zero. All denominators are nonzero
for $n\geq4$. Finally,
$q_0+q_1+q_2=3(111n^2+456n+229)$, so the correction $-4$ supplies
exactly the negative of the inhomogeneous term.
\end{proof}

\begin{corollary}\label{cor:half}
For $b_n=\sum_{k=0}^{\lfloor n/2\rfloor}T(n,k)$, $n\geq0$ (A027179),
put
\begin{align*}
p_0(n)&=-(n+2)(3n-20),&p_1(n)&=-(n+1)(3n-11),\\
p_2(n)&=33n^2-180n-100,&p_3(n)&=15n^2-56n+1,\\
p_4(n)&=2(-51n^2+313n-208),&p_5(n)&=-12(n-1)(n-4),\\
p_6(n)&=8(9n-11)(n-5).
\end{align*}
Then Mathar's seven-term recurrence holds for every $n\geq6$:
\begin{equation}\label{eq:half}
\sum_{j=0}^6p_j(n)b_{n-j}+4(-24n^2+196n-365)=0.
\end{equation}
\end{corollary}

\begin{proof}
Summing the two binomial terms in $T(n,k)$ gives
\[
S_n:=\sum_{k=0}^nT(n,k)
=2(2^{n+2}-2)+2^n-4(n+1)=9\cdot2^n-4n-8.
\]
This is the row-sum formula recorded by Ralf Stephan in 2004 and proved
by Krause in 2026 in A027178~\cite{A027178}.
By row symmetry, the half-row sum is $S_n/2$ when $n$ is odd and
$(S_n+T(n,n/2))/2$ when $n$ is even, as recorded by
Mathar~\cite{A027179}. Define
\[
w_{2m}=d_m=\frac{9m+5}{m+1}c_m,\qquad w_{2m+1}=0\quad(m\geq0).
\]
In even rows, $T(2m,m)=d_m-4$.
Since the indicator of even $n$ is $(1+(-1)^n)/2$, this gives the
all-index closed form
\begin{equation}\label{eq:halfclosed}
b_n=\frac92\,2^n-2n-5-(-1)^n+\frac12w_n\qquad(n\geq0).
\end{equation}
Write $(\mathcal Pf)_n=\sum_{j=0}^6p_j(n)f_{n-j}$. Expanding the
seven coefficients gives
\begin{align*}
\sum_{j=0}^6p_j(n)2^{-j}&=0,&
\sum_{j=0}^6p_j(n)(-1)^j&=0,\\
\sum_{j=0}^6p_j(n)(-2(n-j)-5)&=4(24n^2-196n+365).
\end{align*}
Thus it remains to check $\mathcal Pw=0$.
For even $n\geq6$, only shifts $j=0,2,4,6$ contribute. Dividing by
$c_{n/2}$ and using the ratio for consecutive central binomial
coefficients gives
\[
\frac{(\mathcal Pw)_n}{c_{n/2}}=
\frac{E(n)}{8(n-3)(n-1)},
\]
where the explicit certificate is
\begin{equation}\label{eq:evencert}
\begin{aligned}
E(n)={}&-8(n-3)(n-1)(3n-20)(9n+10)\\
 &+2(n-3)(9n-8)(33n^2-180n-100)\\
 &-n(9n-26)(51n^2-313n+208)\\
 &+n(n-2)(9n-44)(9n-11)=0.
\end{aligned}
\end{equation}
For odd $n\geq7$, only shifts $j=1,3,5$ contribute, and similarly
\[
\frac{(\mathcal Pw)_n}{c_{(n-1)/2}}=\frac{O(n)}{4(n-2)},
\]
with
\begin{equation}\label{eq:oddcert}
\begin{aligned}
O(n)={}&-4(n-2)(3n-11)(9n+1)\\
 &+(9n-17)(15n^2-56n+1)-3(n-1)^2(9n-35)=0.
\end{aligned}
\end{equation}
Both final equalities follow by expanding the displayed products.
For clarity, all terms before clearing denominators can be recovered
from
\[
\frac{w_{n-j}}{c_{(n-\epsilon)/2}}=
\frac{9(n-j)+10}{n-j+2}
\prod_{\ell=0}^{(j-\epsilon)/2-1}
\frac{n-\epsilon-2\ell}{4(n-\epsilon-2\ell-1)},
\quad j\equiv\epsilon\equiv n\pmod2,
\]
where an empty product is $1$. The other shifts give zero.
There is no zero denominator in the stated ranges.
Applying $\mathcal P$ to~\eqref{eq:halfclosed} now proves~\eqref{eq:half}.
\end{proof}

\paragraph{Verification and availability.}
The argument above proves the identities for all stated indices.
The accompanying Python check separately compares the defining convolution
with~\eqref{eq:closed} through $n=400$, checks all five recurrences
through $n=400$, and checks the displayed polynomial certificates exactly.
These finite checks are separate from the all-index proofs above. The
\href{https://github.com/the-omega-institute/trureturing/blob/494f471b70d9470db92b4dbc07a309d78ca59ec7/D5/S1/Recurrence/MatharA027182SelfConvolution.lean}{pinned file \texttt{MatharA027182SelfConvolution.lean}}
(commit \texttt{494f471b70d9}) has the same $n\geq6$ statement as theorem
\texttt{MatharA027182SelfConvolution.result}.
The companion corollaries have the written proofs and exact arithmetic
checks given here.

\paragraph{License.}
This note is licensed under the
\href{https://creativecommons.org/licenses/by/4.0/}{Creative Commons Attribution 4.0 International license (CC BY 4.0)}.

\begin{thebibliography}{9}
\bibitem{A027182}
C. Kimberling, \emph{A027182: a(n) = self-convolution of row n of array
T given by A027170}, The On-Line Encyclopedia of Integer Sequences,
FORMULA section: R.\ J.\ Mathar's conjecture of 6 May 2026.
\url{https://oeis.org/A027182}. Accessed 1 November 2026.

\bibitem{A027170}
C. Kimberling, \emph{A027170: Triangular array T read by rows
(4-diamondization of Pascal's triangle)},
The On-Line Encyclopedia of Integer Sequences.
\url{https://oeis.org/A027170}. Accessed 1 November 2026.

\bibitem{A027171}
C. Kimberling, \emph{A027171: a(n) = A027170(2n, n)},
The On-Line Encyclopedia of Integer Sequences,
FORMULA section: Mathar's conjecture of 6 May 2026 and Krause's closed
form of 10 May 2026. \url{https://oeis.org/A027171}.
Accessed 1 November 2026.

\bibitem{A027174}
C. Kimberling, \emph{A027174: a(n) = A027170(2n, n-2)},
The On-Line Encyclopedia of Integer Sequences,
FORMULA section: Mathar's conjecture of 6 May 2026 and Krause's closed
form of 11 May 2026. \url{https://oeis.org/A027174}.
Accessed 1 November 2026.

\bibitem{A027175}
C. Kimberling, \emph{A027175: a(n) = A027170(2n-1, n-1)},
The On-Line Encyclopedia of Integer Sequences,
FORMULA section: Mathar's conjecture of 6 May 2026 and Krause's closed
form of 11 May 2026. \url{https://oeis.org/A027175}.
Accessed 1 November 2026.

\bibitem{A027179}
C. Kimberling, \emph{A027179: a(n) = T(n,0) + T(n,1) + ... + T(n,[ n/2 ]),
T given by A027170}, The On-Line Encyclopedia of Integer Sequences,
FORMULA section: Mathar's conjecture and parity formulas of 6 May 2026.
\url{https://oeis.org/A027179}. Accessed 1 November 2026.

\bibitem{A027178}
C. Kimberling, \emph{A027178: a(n) = T(n,0) + T(n,1) + ... + T(n,n),
T given by A027170}, The On-Line Encyclopedia of Integer Sequences,
FORMULA section: Stephan's formula of 13 February 2004, with Krause's
proof attribution of 14 May 2026.
\url{https://oeis.org/A027178}. Accessed 1 November 2026.

\bibitem{BLS}
A. Bostan, P. Lairez and B. Salvy, \emph{Multiple binomial sums},
Journal of Symbolic Computation \textbf{80}, part 2 (2017), 351--386.
\href{https://doi.org/10.1016/j.jsc.2016.04.002}{doi:10.1016/j.jsc.2016.04.002}.
Preprint: \href{https://arxiv.org/abs/1510.07487v2}{arXiv:1510.07487v2}.
\end{thebibliography}
\end{document}
